
\subsection{可供借用的一般理論：$M$-估計、混合表示、損失與懲罰的分類，與信賴度標準}\label{subsec:mest}\label{subsec:mix}\label{subsec:mq}\label{subsec:cred}

上一小節的缺口可以一句話寫定：標準 Lee-Carter 的目標函數不是任何分配的
對數概似，故以概似為前提的偏誤修正無從施加。但它仍然是某個目標函數的
極小值，而這一點足以使它落入一個更一般的框架。本小節回顧該框架，
依序為 $M$-估計與 Fisher 一致性、誤差分布的混合表示、
損失與懲罰的一般分類，以及本文結果所要接上的精算實務標準。

\citet{huber1964} 提出 $M$-估計量，把最大概似的想法推廣為「某個目標函數的
極小值」或等價地「某個估計方程的解」：$\hat\theta$ 解
$\sum_i\psi(y_i;\theta)=0$，其中 $\psi$ 不必是對數概似的導數。
該文的動機是位置參數的穩健估計，設定為 $\varepsilon$-污染（絕大多數觀測
來自模型、少數來自別處），目標在於使少數觀測無法主導估計。
此一概念隨後被量化：\citet{hampel1974} 以影響函數刻畫單一觀測的影響力，
\citet{donoho1983} 以崩潰點刻畫可承受的污染比例上限，
\citet{huber2009} 為該領域的標準專著；而穩健性所付的代價為效率，
例如中位數在常態下相對於樣本平均的效率僅 $2/\pi\approx0.64$。
本文推導時所依循的記號取自 \citet[第 1 章]{zoubir2018}：該章把估計量寫成
作用在經驗分布上的泛函 $\hat\theta_N=T(F_N)$（該書式 1.30），
以 $T(F)$ 與真值之差定義漸近偏誤，並於 §1.2、§1.3.1 與 §1.3.2 分別給出
$M$-估計量作為估計方程之解的定義、影響函數與崩潰點。
本文的核心推導並未直接引用該書或以下各書中的現成定理，所借用的是框架與記號。
在應用上，$M$-估計已是訊號處理、電力系統狀態估計與電腦視覺的常規工具
\citep{zoubir2018}，其場景的共同特徵是觀測數遠大於參數數、
且離群值確實是少數。在死亡率文獻中，\citet{santolino2020} 把分位數迴歸
引入 Lee-Carter 模型，屬同一個框架下改動損失函數的作法。

$M$-估計的另一條發展線把「估計方程」本身當成建模的對象，而不必先寫下
概似。\citet{liang1986} 的廣義估計方程即屬此類，
其以工作相關矩陣處理相依資料而不指定聯合分布，在縱貫資料分析中已成常規；
\citet{merlo2022} 進一步把廣義估計方程推廣到帶 Huber 損失的 M-分位數。
這條線近年在因果推論中被推到一般的原則：\citet{chernozhukov2018} 要求所用的
估計方程在擾動參數的真值處對該擾動的一階變動不敏感，
並以樣本切分使該性質在高維擾動下仍成立，其所稱的去偏即是把估計方程本身
重新構造到期望為零。這條線確立了一個看法：一個估計量的性質，
可以完全由其估計方程的期望與變異數讀出，
而不必訴諸該估計量是否為某個概似的最大值。

$M$-估計理論中另有一項比穩健性更為基本的性質，即\textbf{Fisher 一致性}。
一個 $M$-估計量要在真值 $\theta_0$ 處一致，
必要條件是估計方程在真值處的期望為零，
$\Ex_{\theta_0}[\psi(Y;\theta_0)]=0$；否則 $\hat\theta$ 收斂到的是
使該期望為零的另一個點，與 $\theta_0$ 相差一個不隨樣本數消失的量。
此一條件在教科書中通常一行帶過，因為對最大概似而言它自動成立
（分數函數的期望為零），對最小平方而言只要誤差期望為零即成立。
但它之所以自動成立，前提是損失函數作用在未經轉換的反應變數上；
一旦反應變數先經過非線性轉換，$\Ex[\psi]=0$ 就不再是自動的，
而古典文獻對此一情形並未展開。

這條脈絡的標準設定與小人口 Lee-Carter 之間另有幾處落差。
古典 $M$-估計文獻關心的是 $\psi$ 有界所帶來的穩健性，
而小人口 Lee-Carter 所涉的是 $\psi$ 的\textbf{期望}是否為零，
兩者是不同的問題：平方損失的 $\psi(u)=u$ 完全不穩健，
但若 $\Ex[u]=0$ 它仍是一致的；反之一個很穩健的 $\psi$ 若 $\Ex[\psi]\ne0$
則仍不一致。其次，古典設定的觀測數遠大於參數數，而 Lee-Carter 模型每個
$\alpha_x$ 只有 $T$ 個觀測、每個 $\kappa_t$ 只有 $A$ 個，兩者皆為數十的量級。
再者，古典設定的污染比例未知，而零格替代所造成的污染比例則為已知，
見下文。另有一項結構上的落差。古典 $M$-估計的參數進入模型的方式是線性的
（$x_i^{\top}\theta$），而 Lee-Carter 的 $\alpha_x+\beta_x\kappa_t$ 對
$(\beta,\kappa)$ 是\textbf{雙線性}的，識別須靠兩個約束；
這使 $\hat\beta$ 與 $\hat\kappa$ 的漸近理論複雜得多；而年齡水準參數
$\hat\alpha_x$ 不受此影響，它在雙線性部分被估計之前即已由列平均定出，
因此可以完全在單變數 $M$-估計的框架內分析。

污染比例已知這一點，須由混合模型的表示法讀出。有限混合模型把觀測的分布
寫成若干成分分布的加權和，$f(y)=\sum_{j}\pi_j f_j(y)$。
其發展脈絡自 \citet{pearson1894} 以動差法分解螃蟹前額寬的雙峰分布起，
長期受限於計算；\citet{dempster1977} 的 EM 演算法使最大概似估計成為可行，
此後混合模型成為模型式分群的標準工具，\citet{mclachlan2000} 與
\citet{yao2024} 為前後兩代專著，後者涵蓋半參數推廣、穩健建模、
標籤交換與高維情形。本文所借用的即 \citet[§1.2]{yao2024} 的有限混合寫法，
而該書 §1.3 的識別性與 §1.4、§1.5 的最大概似與 EM 演算法皆以各成分有密度
為前提。在應用上，混合模型用於天文、基因體、經濟與金融等領域的異質資料；
在小區域估計中，\citet{delsarto2019} 以 M-分位數迴歸的有限混合處理未觀測
異質性，\citet{denicolo2023} 以 Beta 混合估計地區的不均度指標。

與計數資料最相關的一支是零膨脹模型。\citet{lambert1992} 的零膨脹 Poisson
迴歸把觀測寫成「一個零處的點質量」與「一個 Poisson」的兩成分混合，用以處理
製造業瑕疵計數中零的個數遠多於 Poisson 所能解釋者。該文的形式與零格替代所
造成的混合極為相似，兩者的差別須先釐清。零膨脹模型所處理的是\textbf{結構上
多出來的零}，亦即資料生成機制本身含有一個「必然為零」的狀態，
故零的個數超過 Poisson 的預期。零格替代所造成的混合\textbf{不是}零膨脹：
死亡數就是純粹的 Poisson，零的個數恰如其機率所預期，沒有任何多餘；
該混合是替代規則在取對數之後的變數上造出來的，因為 $\log 0=-\infty$
不可用，分析者必須把所有零格對應到同一個值，
於是轉換後的變數才出現一個點質量。換言之，零膨脹模型的混合來自
\textbf{資料}，零格替代的混合來自\textbf{分析步驟}。
這個差別有實質的性質：零膨脹模型須估計膨脹比例，而零格替代的混合權重由
分布本身給定，無須估計；反之零膨脹模型可以不取對數而直接配適，
零格替代的困難則完全來自取對數這個步驟。

混合模型與穩健統計的連結早已建立：Huber 的 $\varepsilon$-污染模型
$(1-\varepsilon)F+\varepsilon H$ 就是一個兩成分混合。但兩個文獻的處理方式
相反。混合模型把 $\pi_j$ 與 $f_j$ 都當成待估的參數，以 EM 求解；
穩健統計則把 $\varepsilon$ 與 $H$ 當成未知的滋擾，設計對其不敏感的估計量。
穩健估計量之所以必須付出效率代價，正是因為 $\varepsilon$ 未知：
它必須針對最壞的情形設計。
這條脈絡與小人口 Lee-Carter 之間的落差正在此處，而該落差有一半是有利的。
零格替代使每一格的觀測成為一個兩成分混合，而其混合權重是\textbf{已知的}，
它就是死亡數為零的機率，拿到資料即可估計；
既然如此，既不必以 EM 去估它，也不必為不知道它而承擔該項效率損失。
另一半則不利：該混合的其中一個成分是\textbf{退化的}（一個點質量），
而有限混合理論通常假設各成分有密度，故識別性、標籤交換等標準議題在此
不適用，EM 亦無從施加。
另有兩項性質由混合權重已知而來。其一，\citet{donoho1983} 的 $50\%$ 上界
可直接換算為期望死亡數的門檻，不必以模擬估計；其二，零格替代所造成的污染
比例在幼年組可超過七成，故穩健統計「污染為少數」這個前提在小人口死亡資料上
並不成立。

最後一個可借用的結果是分類。把損失函數、權重與懲罰整合為同一個目標函數的
作法，在小區域估計中發展得最完整。
\citet{koenker1978} 以檢查函數 $\rho_\tau$ 取代平方損失，
使迴歸得以刻畫條件分位數；\citet{koenker2005} 為標準參考，
其應用場景包括勞動經濟學的薪資分布與兒童生長參考曲線
\citep{wei2006}，共同特徵是每個共變數水準上有大量觀測。
該書第 5 章（頁 151--172）給出加權的最適形式，其中定理 5.1 指出分位數迴歸的
權重應與所關心分位處的局部密度成比例，而非與觀測標準差的倒數成比例；
第 8 章（頁 250--292）另討論離散反應與高崩潰點的替代方案，
並說明離散反應使該書標準設定下的漸近結果不再適用的原因。
\citet{breckling1988} 指出，若把 Huber 損失乘上一個依殘差正負
而異的非對稱權重，即得一個由 $(\tau,c)$ 索引的損失函數族 M-分位數，
其兩個極限分別是檢查函數（$c\to0$）與非對稱最小平方
（$c\to\infty$，\citealp{newey1987}）。

M-分位數在小區域估計中受青睞的理由如下。
傳統的小區域估計以 \citet{fay1979} 的隨機效果模型為主，
須假設區域效果的分布（通常為常態），其均方誤差的估計則由
\citet{prasad1990} 給出；兩者的收縮量由變異成分決定而無須指定懲罰強度，
但其推導假設直接估計量不偏、僅帶抽樣噪音。
\citet{chambers2006} 指出 M-分位數可直接給出各區域
自己的迴歸係數而不必作此分布假設，因而對分布誤設較不敏感；
該法的應用多為家戶調查資料的地方層級所得與剝奪指標估計。
這條線的後續發展相當密集：\citet{bianchi2018} 補上推論理論，
\citet{ranalli2023} 加入 Lasso 與 Elastic Net 懲罰並應用於空氣品質，
\citet{hu2021} 以損失的冪次為刻度提出 $L_p$-分位數並處理高維懲罰，
\citet{merlo2022} 處理多變量相依，\citet{dawber2022} 處理離散標的，
\citet{bugallo2026} 針對均方誤差估計提出偏誤校正，
\citet{bugallo2025} 另證明區域別係數的一致性並處理殘差診斷，
\citet{lahiri2023} 則把巢式誤差迴歸推廣到同時容許區域別係數與區域別變異
成分的情形；\citet{jiang2026} 為這一整支穩健小區域估計的近期專著。

與此平行的是正則化一脈，其想法在小區域估計之外獨立發展。
\citet{stein1956} 指出高維常態平均數的樣本平均不可容許，
\citet{james1961} 給出具體的收縮估計量，
\citet{hoerl1970} 的 Ridge 與 \citet{tibshirani1996} 的 LASSO 則把收縮
寫成在目標函數上加懲罰項，\citet{zou2005} 與 \citet{zou2006} 再加以擴充。
\citet{fay1979} 的經驗貝氏收縮即是同一想法的區域層級版本。
這一脈的懲罰強度通常以交叉驗證選定，其標準論述見 \citet{buhlmann2011} 與
\citet{hastie2015}。把穩健損失與懲罰接起來的整理見 \citet{zoubir2018} 第 3 章，
其中並給出迴歸係數與尺度的聯合懲罰式 $M$-估計。
綜合之，這條脈絡的一般目標函數為
\begin{equation}\label{eq:general}
\min_{\theta}\ \sum_{i} w_i\,\rho_{\tau,c}(u_i(\theta))\;+\;\lambda P(\theta),
\end{equation}
其中含三項設定：損失形狀 $(\tau,c)$、權重 $w$、懲罰 $(\lambda,P)$。

這條脈絡的標準設定與小人口 Lee-Carter 之間的落差亦有數處。
其應用場景是單位層級的調查資料，每個小區域有多個受訪單位；
而死亡率資料是聚合的計數，每格只有一個觀測。
其次，三項設定都需要調參（$\tau$、$c$、$\lambda$），
而死亡率資料既無真值、亦無法以交叉驗證選參，蓋因切分年齡--年份格會破壞
雙線性結構所依賴的跨格借力，被留置的格無法由其餘的格預測。
再者，這一脈的漸近理論均假設解釋變數已觀測，而 Lee-Carter 的 $\kappa_t$ 是
待估的潛在變數而非觀測值，故上述理論不能直接套用。
至於式~\eqref{eq:general} 本身，它把標準 Lee-Carter 納為一個角落
（$\tau=0.5$、$c=\infty$、$w\equiv1$、$\lambda=0$），
因而可以用來辨認各種補救各自改動了其中哪一項設定：
第\ref{sec:ref}所回顧的修勻改動的是資料，而本小節所列的各種方法改動的是
損失形狀、權重或懲罰。

最後一條脈絡不是統計推導所依的理論，而是精算的實務標準，
亦即本文結果所要接上的判準。保險業判斷一組資料是否足以單獨用於費率計算，
其依循的即是信賴度理論。全信賴度標準把資料量換算為期望索賠次數，
並要求該次數高到使估計量的相對誤差以指定的機率落在指定的範圍內；
在 $k=5\%$、$p=90\%$ 的慣例下，該門檻約為 $1{,}082$ 件
（\citealp{klugman2012}，第 17 章）。部分信賴度則以該門檻的平方根比例在
目標資料與參考資料之間插值，其形式與第\ref{sec:ref}的 partial SMR 同型。
此一標準有一項性質值得指出：\textbf{它控制的是估計的變異，
而其推導假設估計量無偏}。全信賴度的計算只用到估計量的變異數與常態近似，
偏誤不出現在任何一條式子裡。這在索賠次數的設定下是合理的，
因為該處的估計量就是樣本平均；但在小人口 Lee-Carter 的設定下並不成立，
蓋因其估計量在零格存在時即非無偏，而偏誤的量與資料量的關係並不隨資料量
增加而自動可忽略。就本文所查，信賴度文獻中並無與全信賴度標準對應的
\textbf{偏誤}門檻，這條脈絡所留下的缺口即在此處。
